% 1_restatement.tex — 问题重述

\subsection{问题背景}

烟幕干扰弹通过在目标前方特定空域形成烟幕或气溶胶云团来遮蔽来袭武器的探测系统，具有成本低、效费比高的优点。在现代化战争中，利用无人机挂载烟幕干扰弹在来袭武器和保护目标之间形成有效遮蔽，是一种重要的防御手段。

\subsection{场景描述}

考虑以下作战场景：
\begin{itemize}
    \item 假目标位于坐标原点$O(0,0,0)$，为半径7 m、高10 m的圆柱形固定目标；
    \item 真目标下底面圆心位于$(0,200,0)$，同样为半径7 m、高10 m的圆柱形目标；
    \item 来袭导弹为空地导弹，飞行速度300 m/s，飞行方向直指假目标；
    \item 警戒雷达发现来袭导弹时，3枚导弹M1、M2、M3分别位于$(20000,0,2000)$、$(19000,600,2100)$、$(18000,-600,1900)$（单位：m）；
    \item 5架无人机FY1$\sim$FY5分别位于$(17800,0,1800)$、$(12000,1400,1400)$、$(6000,-3000,700)$、$(11000,2000,1800)$、$(13000,-2000,1300)$。
\end{itemize}

无人机受领任务后可瞬时调整飞行方向，以70$\sim$140 m/s的速度等高度匀速直线飞行，每架无人机的航向和速度一旦确定不再调整。烟幕干扰弹脱离无人机后在重力作用下运动，起爆后瞬时形成半径10 m的球状烟幕云团，以3 m/s速度匀速下沉，起爆后20 s内可为目标提供有效遮蔽。每架无人机投放两枚烟幕干扰弹至少间隔1 s。

\subsection{需要解决的问题}

\textbf{问题1：}利用FY1以120 m/s速度朝向假目标飞行，受领任务1.5 s后投放1枚烟幕干扰弹，间隔3.6 s后起爆，计算对M1的有效遮蔽时长。

\textbf{问题2：}利用FY1投放1枚烟幕干扰弹对M1进行干扰，确定最优的飞行方向、飞行速度、投放点和起爆点，使遮蔽时间尽可能长。

\textbf{问题3：}利用FY1投放3枚烟幕干扰弹对M1进行干扰，给出投放策略并输出到result1.xlsx。

\textbf{问题4：}利用FY1、FY2、FY3各投放1枚烟幕干扰弹对M1进行干扰，给出投放策略并输出到result2.xlsx。

\textbf{问题5：}利用5架无人机（每架至多投放3枚弹）对M1、M2、M3进行干扰，给出投放策略并输出到result3.xlsx。
